\documentclass[a4paper,11pt]{article}
\usepackage{petritikz}
\usepackage{amsmath}
\title{\textbf{Reseau de Petri}}
\author{PetriTikz}
\date{\today}
\begin{document}
\maketitle

\begin{center}
% PetriTikz | layout=circle
% Pre  = [[1, 0, 0, 0], [0, 1, 1, 0]]
% Post = [[0, 1, 1, 0], [1, 0, 0, 1]]
\begin{tikzpicture}[petrinet]

  % Places  (positions calculees par Python depuis Pre/Post)
  \node[place,label=above:libre] (p1) at (0.0,3.5) {};
  \petritokens{p1}{1}
  \node[place,label=above:produit] (p2) at (-3.5,0.0) {};
  \node[place,label=above:vide] (p3) at (-0.0,-3.5) {};
  \petritokens{p3}{1}
  \node[place,label=above:plein] (p4) at (3.5,-0.0) {};

  % Transitions
  \node[enabled transition,label=below:produire] (t1) at (-2.0,0.0) {};
  \node[transition,label=below:consommer] (t2) at (2.0,-0.0) {};

  % Arcs Pre : P -> T  (depuis net.pre)
  \draw[arc] (p1) -- (t1);
  \draw[arc] (p2) -- (t2);
  \draw[arc] (p3) -- (t2);

  % Arcs Post : T -> P  (depuis net.post)
  \draw[arc] (t1) -- (p2);
  \draw[arc] (t1) -- (p3);
  \draw[arc] (t2) -- (p1);
  \draw[arc] (t2) -- (p4);

\end{tikzpicture}
\end{center}

\[ \text{Pre}=\begin{pmatrix}
     \\1 & 0 & 0 & 0
    0 & 1 & 1 & 0
\end{pmatrix}\quad\text{Post}=\begin{pmatrix}
     \\0 & 1 & 1 & 0
    1 & 0 & 0 & 1
\end{pmatrix}\quadC=\begin{pmatrix}
     \\-1 & 1 & 1 & 0
    1 & -1 & -1 & 1
\end{pmatrix}\]

\[ M_0=(1,0,1,0) \qquad \mathcal{T}_{\text{act}}=\{t1\} \]

\end{document}
